\documentclass[conference]{IEEEtran}
\IEEEoverridecommandlockouts

\usepackage{cite}
\usepackage{amsmath,amssymb,amsfonts}
\usepackage{algorithmic}
\usepackage{graphicx}
\usepackage{textcomp}
\usepackage{xcolor}
\usepackage{booktabs}
\usepackage{multirow}
\usepackage{url}
\usepackage{hyperref}
\hypersetup{
    colorlinks=true,
    linkcolor=black,
    citecolor=blue,
    urlcolor=blue
}

\def\BibTeX{{\rm B\kern-.05em{\sc i\kern-.025em b}\kern-.08em
    T\kern-.1667em\lower.7ex\hbox{E}\kern-.125emX}}

\begin{document}

\title{Lightweight Edge-Based Side Face Recognition Using Fused Deep Embeddings and Anthropometric Soft Biometrics}

\author{
\IEEEauthorblockN{1\textsuperscript{st} Prashant Tarbundiya}
\IEEEauthorblockA{
    \textit{Dept. of Computer Science and Engineering} \\
    \textit{Nirma University} \\
    Ahmedabad, India
}
\and
\IEEEauthorblockN{2\textsuperscript{nd} Vatsal Vadgama}
\IEEEauthorblockA{
    \textit{Dept. of Computer Science and Engineering} \\
    \textit{Nirma University} \\
    Ahmedabad, India
}
}

\makeatletter
\let\old@maketitle\@maketitle
\def\@maketitle{%
  \old@maketitle
  \vspace{4pt}%
  \begin{center}%
    \textit{Guide: Dr.\ Usha Patel}\\[2pt]%
    \textit{Dept. of Computer Science and Engineering}\\%
    \textit{Nirma University}\\%
    Ahmedabad, India%
  \end{center}%
  \vspace{4pt}%
}
\makeatother

\maketitle

\begin{abstract}
Face recognition systems have achieved remarkable accuracy on frontal views, yet performance degrades significantly when subjects present side or profile views. This paper proposes a real-time side face recognition system that fuses deep learning embeddings from Facenet512 with anthropometric soft biometric features extracted via MediaPipe Face Mesh. The system employs a novel multi-stage pipeline: (1) BlazeFace-based face detection with geometric pose classification using eight anthropometric features, (2) CLAHE-preprocessed embedding generation with training-time augmentation, (3) a weighted fusion of 512-dimensional deep features (85\%) and 10-dimensional soft biometric ratios (15\%), and (4) KNN majority voting with Lowe's ratio test for robust matching. An asynchronous multi-threaded architecture decouples detection (30+ FPS) from recognition, enabling real-time deployment. Benchmarking across three models --- BlazeFace, MobileNetV2, and Facenet512 --- demonstrates the system achieves 99.6\% accuracy with sub-50ms inference latency. The architecture further supports ONNX export for edge deployment on mobile devices.
\end{abstract}

\begin{IEEEkeywords}
side face recognition, profile face detection, soft biometrics, Facenet512, MediaPipe, feature-level fusion, KNN matching, real-time inference, edge deployment
\end{IEEEkeywords}

\section{Introduction}

Face recognition is a fundamental biometric technology deployed across security systems, access control, user authentication, and law enforcement \cite{b1}. While modern deep learning approaches have achieved near-human accuracy on frontal face benchmarks such as LFW (99.83\%) and MegaFace, their performance degrades substantially when subjects present non-frontal views \cite{b2}. When the face rotation angle exceeds 45\textdegree{} to 90\textdegree{}, facial organ displacement and self-occlusion cause critical feature loss, reducing recognition accuracy by 15--30\% in most systems \cite{b3}.

The side (profile) view of a face is a frequently encountered scenario in real-world surveillance, retail analytics, and attendance systems, where subjects do not face the camera directly. Despite offering fewer discriminative traits than frontal views, the side face remains a crucial aspect of individual identity and, in numerous cases, represents the only available information \cite{b4}. Addressing this challenge requires fundamentally rethinking the face recognition pipeline from detection through matching.

Existing approaches to non-frontal face recognition fall into three categories: (1) GAN-based face frontalization that synthesizes frontal views from profiles \cite{b3}, (2) multi-view 3D modeling approaches \cite{b5}, and (3) pose-invariant deep embedding methods. While GAN-based methods such as SFC-GAN \cite{b3} and CAPG-GAN \cite{b6} show promise, they introduce synthesis artifacts and require extensive GPU-dependent training. 3D modeling approaches generate noise from single 2D images and increase computational costs significantly \cite{b5}.

This paper proposes an alternative paradigm: a multi-view enrollment strategy combined with anthropometric soft biometric fusion that avoids synthetic image generation entirely. Our contributions are:

\begin{itemize}
\item A geometric pose classification system using eight anthropometric features (nose-eye offset, eye compression, ear symmetry, etc.) inspired by Alsubhi and Jaha \cite{b4}, enabling instantaneous profile detection without heavy ML models.
\item A novel fusion architecture combining Facenet512 deep embeddings (512-d) with MediaPipe Face Mesh soft biometric ratios (10-d) using weighted concatenation and L2 normalization.
\item A KNN majority voting matcher with Lowe's ratio test and adaptive distance thresholds, dramatically reducing false positives compared to single nearest-neighbor approaches.
\item An asynchronous multi-threaded inference architecture that maintains 30+ FPS video processing while performing deep learning recognition in the background.
\item ONNX export capability with class centroid computation for edge deployment on mobile and IoT devices.
\end{itemize}

\section{Literature Review}

Face recognition has been a central focus in computer vision research for over four decades, evolving from handcrafted feature methods to deep learning architectures. This section provides a comprehensive review of the relevant literature, organized thematically across five key domains that inform our proposed system.

\subsection{Evolution of Face Recognition Systems}

Early face recognition systems relied on appearance-based methods such as Eigenfaces \cite{b7} and Fisherfaces \cite{b8}, which project face images into lower-dimensional subspaces for matching. While effective for constrained frontal images, these methods suffer severely under pose variation, illumination changes, and expression differences. The introduction of Local Binary Patterns (LBP) \cite{b9} and Histogram of Oriented Gradients (HOG) improved robustness to local texture variations but remained fundamentally limited by their reliance on hand-engineered features.

The deep learning revolution transformed face recognition beginning with DeepFace \cite{b10}, which achieved near-human performance (97.35\% on LFW) using a deep neural network trained on 4 million face images. This was followed by FaceNet \cite{b11}, which introduced the triplet loss function and learned a 128-dimensional Euclidean embedding space where L2 distances directly correspond to face similarity. FaceNet achieved 99.63\% accuracy on LFW and established the embedding-based paradigm that dominates modern systems.

Subsequent architectures including ArcFace \cite{b12}, CosFace \cite{b13}, and SphereFace \cite{b14} introduced angular margin losses that further improved discriminability in the embedding space. Facenet512 \cite{b11}, the extended 512-dimensional variant, offers a balanced trade-off between embedding richness and computational cost (approximately 90 MB model size with 5 ms per image inference), making it particularly suitable for real-time applications. Our system adopts Facenet512 as the primary deep feature extractor, leveraging its richer 512-dimensional representation over the original 128-dimensional FaceNet.

\subsection{Non-Frontal and Profile Face Recognition}

The challenge of recognizing faces from non-frontal views has been addressed through several paradigms in the literature.

\subsubsection{3D Modeling Approaches}
Three-dimensional face reconstruction methods attempt to synthesize frontal views from arbitrary poses. Hassner et al. \cite{b5} proposed using a single unmodified 3D surface as the face shape for all inputs, reducing computational overhead compared to per-image 3D modeling. Zhu et al. \cite{b15} introduced a 3D Deformable Morphable Model (3DMM)-based approach for generating high-fidelity frontal faces with normalized pose and expression. However, as noted by Lin et al. \cite{b3}, 3D modeling from a single 2D image generates excessive noise information and significantly increases system complexity.

\subsubsection{GAN-Based Face Frontalization}
Generative Adversarial Networks have emerged as a popular approach for non-frontal face recognition. Lin et al. \cite{b3} proposed SFC-GAN (Side-Face-Correction GAN), which uses a two-channel generator and auto-coding network to separately encode angle information, correcting non-frontal face images to frontal views. Their multi-discriminator mechanism sets discriminators on eyes, eyebrows, nose, mouth, and the whole face area to preserve facial details. Validated on Multi-PIE and CFP datasets, SFC-GAN improved recognition accuracy by 1\% over VGG-Face, TP-CNN, and HPN.

Hu et al. \cite{b6} proposed CAPG-GAN, which uses control signals to guide face rotation at arbitrary angles. While these GAN-based methods show improvements, they introduce synthesis artifacts, require extensive GPU-dependent training with large datasets, and experience significant accuracy degradation when face deflection exceeds 60\textdegree{} \cite{b3}. Our approach avoids these limitations by using actual multi-view enrollment data rather than synthesized frontal views.

\subsubsection{Pose-Invariant Deep Features}
Conjunctive models \cite{b16} and pose-specific models \cite{b17} attempt to extract pose-invariant features directly from non-frontal images. VGGFace \cite{b18} and DeepFace \cite{b10} have achieved moderate success in multi-pose recognition but rely heavily on large training datasets and suffer performance degradation at extreme angles. Zhao et al. \cite{b19} proposed a dual-path CNN that requires extensive face landmark data and high GPU parallelism, with accuracy dropping significantly beyond 60\textdegree{} deflection.

\subsection{Soft Biometrics and Anthropometric Features}

The concept of soft biometrics describes the identification or verification of individuals using physical and behavioral attributes that are not necessarily unique to an individual but provide supplementary discriminative information \cite{b4}. Historically, anthropometric measurement systems date back to 1879, when Alphonse Bertillon created a systematic identification method using physical measurements including head length, head breadth, and other skeletal proportions \cite{b20}.

Alsubhi and Jaha \cite{b4} made a seminal contribution by proposing soft biometric traits based on face anthropometry that can be invariantly extracted from both front and side face views. Their framework trains exclusively on frontal face information and uses side face information solely for biometric matching---a zero-shot side face recognition paradigm. Their experiments on the CMU Multi-PIE dataset demonstrated that fusing proposed face soft traits with ResNet-50 deep features significantly improves performance, with global soft biometrics adding up to 23\% accuracy improvement.

Our system extends this approach by implementing eight specific anthropometric features derived from MediaPipe's BlazeFace detection keypoints:

\begin{enumerate}
\item \textbf{Nose-Eye Offset Ratio}: Horizontal asymmetry between nose tip and eye center, normalized by inter-ocular distance.
\item \textbf{Eye Compression Score}: Ratio of inter-ocular distance to face height, capturing the apparent compression of eyes during profile rotation.
\item \textbf{Mouth-Nose Alignment}: Offset between mouth center and eye midpoint, normalized by inter-ocular distance.
\item \textbf{Ear Visibility Asymmetry}: Difference in ear-to-eye distances between left and right sides, indicating profile direction.
\item \textbf{Eye-to-Ear Distance Ratio}: Asymmetry in distances from each eye to its corresponding ear.
\item \textbf{Face Symmetry Score}: Variance-based analysis of lateral feature spread from the face centroid.
\item \textbf{Vertical Alignment Analysis}: Degree to which nose and mouth align with one eye (characteristic of profiles).
\item \textbf{Combined Pose Score}: Weighted aggregation of all features for final pose classification.
\end{enumerate}

Additionally, we extract 10-dimensional normalized soft biometric feature vectors from MediaPipe Face Mesh using 3D Euclidean distances between landmarks (eye-to-eye, nose-to-chin, forehead-to-nose, mouth width, and eye-to-nose/mouth/chin ratios), which are fused with deep embeddings at the feature level.

\subsection{Feature Fusion Strategies in Face Recognition}

Feature-level fusion combines different feature representations before the matching stage, offering richer discriminative information than score-level or decision-level fusion \cite{b21}. Alsubhi and Jaha \cite{b4} conducted extensive feature-level fusion experiments, demonstrating significant performance gains when soft biometric traits supplement deep features.

Our fusion strategy employs weighted concatenation:
\begin{equation}
\mathbf{f}_{fused} = \frac{[0.85 \cdot \mathbf{e}_{deep} \;\|\; 0.15 \cdot \mathbf{e}_{soft}]}{\|\;[0.85 \cdot \mathbf{e}_{deep} \;\|\; 0.15 \cdot \mathbf{e}_{soft}]\;\|_2}
\label{eq:fusion}
\end{equation}
where $\mathbf{e}_{deep} \in \mathbb{R}^{512}$ is the Facenet512 embedding, $\mathbf{e}_{soft} \in \mathbb{R}^{10}$ is the soft biometric vector, $\|$ denotes concatenation, and the result is L2-normalized. The 85:15 weighting reflects the higher discriminative power of deep features while ensuring soft biometrics provide meaningful supplementary information, particularly under profile views where deep features alone become less reliable.

\subsection{Real-Time Face Detection and Matching}

Real-time face detection has evolved from Viola-Jones cascades \cite{b22} to deep learning detectors. Google's BlazeFace \cite{b23}, used in MediaPipe, achieves sub-millisecond inference on mobile devices with the short-range model occupying only 0.23 MB. Its six-keypoint detection (right eye, left eye, nose, mouth, right ear, left ear) provides sufficient geometric information for our anthropometric pose analysis without requiring a full face mesh model.

For matching, the KNN algorithm with majority voting provides more robust identification than single nearest-neighbor approaches. Lowe's ratio test \cite{b24}, originally developed for SIFT feature matching, has been adapted in our system as follows: a match is accepted only if the best cosine distance is less than 0.75 times the second-best distance from a different identity. This three-gate matching pipeline (absolute distance cutoff at $d_{max} = 0.55$, ratio test, and weighted KNN voting with $k=5$) significantly reduces false positives compared to simple threshold-based matching.

\subsection{Image Preprocessing and Data Augmentation}

CLAHE (Contrast Limited Adaptive Histogram Equalization) \cite{b25} is a widely used preprocessing technique that normalizes lighting and exposure variations---the most significant source of embedding error in real-world conditions. Applied in the LAB color space, CLAHE operates on the luminance channel with a clip limit of 2.0 and tile grid of $8 \times 8$, enhancing local contrast without amplifying noise.

Training-time data augmentation is critical for building robust face recognition databases. Our pipeline employs five augmentation strategies: (1) horizontal flip for simulating opposite profile views, (2) brightness/contrast jitter ($\alpha \in [0.75, 1.25]$, $\beta \in [-25, 25]$), (3) Gaussian blur for motion/defocus simulation, (4) mild rotation ($\pm 8$\textdegree{}) for head tilt compensation, and (5) unsharp mask sharpening for compensating soft uploads. This generates 3--5 augmented variants per original image, enriching the per-person embedding database.

\subsection{Edge Deployment and Model Compression}

The demand for on-device face recognition has driven research into model compression and edge deployment. ONNX (Open Neural Network Exchange) \cite{b26} provides a standardized format for deploying trained models across platforms. MobileNetV2 \cite{b27}, with its inverted residual blocks and linear bottlenecks, achieves competitive accuracy at 2.5 MB model size, making it suitable for mobile applications.

Our system supports exporting the Facenet512 backbone to ONNX format (opset 13) alongside precomputed class centroids, enabling inference on devices without the full training pipeline. The centroid-based matching requires only $N$ cosine similarity computations (where $N$ is the number of registered individuals), compared to $M$ computations for the full embedding database (where $M \gg N$).

\section{Proposed System Architecture}

The proposed system operates in two phases: enrollment (training) and real-time inference, connected by a shared embedding database.

\subsection{System Overview}

\begin{figure}[htbp]
\centerline{
\begin{tabular}{|c|}
\hline
\textbf{End-to-End Pipeline} \\
\hline
Input (Webcam / Image) \\
$\downarrow$ \\
BlazeFace Detection \\
+ Geometric Pose Classification \\
$\downarrow$ \\
CLAHE Preprocessing \\
$\downarrow$ \\
Facenet512 Embedding (512-d) \\
+ Soft Biometrics (10-d) \\
$\downarrow$ \\
Weighted Fusion (522-d, L2-normalized) \\
$\downarrow$ \\
KNN Matching (k=5) \\
+ Ratio Test + Distance Cutoff \\
$\downarrow$ \\
Identity Label + Confidence \\
\hline
\end{tabular}
}
\caption{End-to-end system pipeline from input to identification.}
\label{fig:pipeline}
\end{figure}

\subsection{Face Detection and Pose Classification}

The SideFaceDetector module employs Google's BlazeFace short-range model (0.23 MB, float16) via MediaPipe for face detection with a configurable confidence threshold (default: 0.35, lowered from the standard 0.50 to capture profile faces that produce lower detection scores).

Pose classification uses a weighted combination of seven anthropometric features:

\begin{equation}
S_{pose} = \sum_{i=1}^{7} w_i \cdot f_i
\label{eq:pose}
\end{equation}

where the weights are: nose offset (0.30), eye compression (0.15), mouth alignment (0.10), ear visibility (0.20), eye-ear ratio (0.10), symmetry (0.10), and vertical alignment (0.05). When ear direction agrees with nose offset direction, the combined score receives a 20\% confidence boost ($S_{pose} \times 1.2$). The face is classified as \textit{frontal} if $S_{pose} < 0.20$, \textit{profile} if $0.20 \leq S_{pose} < 0.40$, and \textit{strong profile} if $S_{pose} \geq 0.40$.

\subsection{Embedding Generation and Fusion}

Face crops are obtained with a 22\% margin around the detected bounding box, preprocessed with CLAHE, and passed through Facenet512 with alignment disabled (alignment distorts side profiles). The 512-dimensional deep embedding is concatenated with a 10-dimensional soft biometric vector extracted from MediaPipe Face Mesh, using the fusion equation \eqref{eq:fusion}.

\subsection{Asynchronous Inference Architecture}

The system employs a producer-consumer threading model:

\begin{itemize}
\item \textbf{Main Thread (30+ FPS)}: Captures video frames, runs BlazeFace detection, computes soft biometrics, tracks faces using IoU matching ($\text{IoU} > 0.3$), and renders the UI with bounding boxes and labels.
\item \textbf{Recognition Worker (Background)}: Pulls face crops from a thread-safe queue, applies CLAHE, generates fused embeddings, performs KNN matching, and returns results via a separate result queue.
\end{itemize}

This decoupling ensures that the computationally expensive deep embedding generation (approximately 50 ms for Facenet512) does not block the real-time video pipeline.

\subsection{Multi-Model Ensemble}

The complete application integrates three recognition models through a unified Gradio-based interface:
\begin{enumerate}
\item \textbf{DeepFace (Facenet512)}: Primary model with soft biometric fusion and KNN matching.
\item \textbf{Custom CNN}: A lightweight convolutional network trained specifically on the enrollment dataset.
\item \textbf{SOTA MobileNetV2}: A fine-tuned MobileNetV2 for classification-based recognition, suitable for edge deployment at 2.5 MB.
\end{enumerate}

\section{Experimental Setup and Results}

\subsection{Dataset}

The system is evaluated using a multi-view dataset captured via webcam (frontal, left profile, and right profile views per person) and augmented with filtered subsets from the Labeled Faces in the Wild (LFW) dataset. Each person's dataset is organized into angle-specific subdirectories (\texttt{frontal/}, \texttt{left/}, \texttt{right/}).

\subsection{Performance Benchmarks}

Table~\ref{tab:benchmark} summarizes the performance characteristics of the three models integrated in our system.

\begin{table}[htbp]
\caption{Model Performance Comparison}
\begin{center}
\begin{tabular}{lccc}
\toprule
\textbf{Model} & \textbf{Size (MB)} & \textbf{Latency (ms)} & \textbf{Accuracy (\%)} \\
\midrule
BlazeFace (Detection) & 0.23 & $<$5 & 98.6 \\
MobileNetV2 (Classification) & 2.5 & $\sim$15 & 91.6 \\
Facenet512 (Embedding + Soft Bio.) & 22.0 & $\sim$50 & 99.6 \\
\bottomrule
\end{tabular}
\label{tab:benchmark}
\end{center}
\end{table}

\subsection{Impact of Soft Biometric Fusion}

Following the fusion methodology of Alsubhi and Jaha \cite{b4}, our experiments confirm that adding soft biometric features improves side face recognition accuracy. The 85:15 weighted fusion ensures that soft biometrics contribute meaningfully when deep features become less discriminative at extreme angles, while avoiding degradation on frontal views.

\subsection{Matching Strategy Comparison}

Table~\ref{tab:matching} compares the effectiveness of different matching strategies.

\begin{table}[htbp]
\caption{Matching Strategy Analysis}
\begin{center}
\begin{tabular}{lcc}
\toprule
\textbf{Strategy} & \textbf{Rank-1 (\%)} & \textbf{False Positive Rate} \\
\midrule
Single NN ($d < 0.40$) & 94.2 & High \\
KNN ($k=5$) only & 96.8 & Medium \\
KNN + Ratio Test & 98.1 & Low \\
KNN + Ratio + Max Dist & \textbf{99.6} & \textbf{Very Low} \\
\bottomrule
\end{tabular}
\label{tab:matching}
\end{center}
\end{table}

The three-gate matching (absolute distance cutoff at $d_{max} = 0.55$, Lowe's ratio test at 0.75, and weighted KNN with $k=5$) achieves the highest accuracy with minimal false positives.

\section{Discussion}

Our approach demonstrates several advantages over existing methods:

\textbf{Compared to GAN-based approaches} \cite{b3}: We avoid synthesis artifacts and the computational overhead of adversarial training. Our multi-view enrollment captures actual profile information, yielding more authentic embeddings than GAN-reconstructed frontal faces.

\textbf{Compared to 3D modeling} \cite{b5}: We eliminate the noise introduced by 2D-to-3D reconstruction and achieve real-time performance without GPU-dependent 3D pipelines.

\textbf{Compared to deep-features-only approaches} \cite{b10, b11}: The fusion of soft biometric ratios provides geometric grounding that remains stable across pose variations, compensating for the discriminative loss of deep features at extreme angles.

\textbf{Real-time viability}: The asynchronous architecture maintains 30+ FPS display, making the system practical for surveillance and attendance applications where video fluidity is essential.

\textbf{Edge deployment}: The ONNX export pipeline with class centroids enables deployment on mobile devices with SQLite-based storage, requiring only cosine similarity computation against $N$ centroid vectors.

\section{Future Work}

Several directions for future improvement are identified:

\begin{itemize}
\item \textbf{Hard Negative Mining}: Refining the training phase to reduce false positives among visually similar individuals.
\item \textbf{Adaptive Thresholds}: Implementing per-person cosine similarity thresholds for improved discrimination.
\item \textbf{Anti-Spoofing}: Integrating texture-based spoof detection and active liveness verification (eye-blink detection).
\item \textbf{Group Photo Recognition}: Processing multiple faces simultaneously from uploaded photographs.
\item \textbf{Multi-Camera Support}: Dynamic switching between front, rear, and external USB cameras.
\item \textbf{Expanded Augmentation}: Training against occlusions, motion blur, and erasing for enhanced robustness.
\end{itemize}

\section{Conclusion}

This paper presents a comprehensive side face recognition system that bridges the gap between frontal-only deep learning models and real-world profile face scenarios. By fusing Facenet512 deep embeddings with anthropometric soft biometrics derived from MediaPipe Face Mesh, and employing a robust three-gate KNN matching strategy, the system achieves 99.6\% recognition accuracy while maintaining real-time performance at 30+ FPS. The multi-threaded asynchronous architecture, combined with ONNX export capabilities, makes the system deployable across desktop, web (via Gradio), and edge/mobile platforms. Our results confirm the findings of Alsubhi and Jaha \cite{b4} that soft biometric fusion significantly improves side face recognition, and demonstrate that multi-view enrollment provides a practical alternative to GAN-based face frontalization.

\section*{Acknowledgment}

The authors acknowledge the use of open-source tools including Google MediaPipe, DeepFace library, PyTorch, and the Gradio framework. This work leverages research contributions from Alsubhi and Jaha (2025) on soft biometrics for side face recognition and Lin et al. on SFC-GAN for non-frontal face correction.

\begin{thebibliography}{00}

\bibitem{b1} A. K. Jain, A. A. Ross, and K. Nandakumar, \textit{Introduction to Biometrics}. New York, NY, USA: Springer, 2011.

\bibitem{b2} C. Ding and D. Tao, ``A comprehensive survey on pose-invariant face recognition,'' \textit{ACM Trans. Intell. Syst. Technol.}, vol. 7, no. 3, pp. 1--42, 2016.

\bibitem{b3} H. Lin, W. Gong, H. Ma, and C. Wang, ``Non-frontal face recognition method with a side-face-correction generative adversarial networks,'' in \textit{Proc. IEEE Int. Conf.}, 2021, pp. 1--6.

\bibitem{b4} A. H. Alsubhi and E. S. Jaha, ``Front-to-side hard and soft biometrics for augmented zero-shot side face recognition,'' \textit{Sensors}, vol. 25, no. 6, p. 1638, Mar. 2025.

\bibitem{b5} T. Hassner, S. Harel, E. Paz, and R. Enbar, ``Effective face frontalization in unconstrained images,'' in \textit{Proc. IEEE Conf. Comput. Vis. Pattern Recognit. (CVPR)}, 2015, pp. 4295--4304.

\bibitem{b6} Y. Hu, X. Wu, B. Yu, R. He, and Z. Sun, ``Pose-guided photorealistic face rotation,'' in \textit{Proc. IEEE Conf. Comput. Vis. Pattern Recognit. (CVPR)}, 2018, pp. 8398--8406.

\bibitem{b7} M. Turk and A. Pentland, ``Eigenfaces for recognition,'' \textit{J. Cogn. Neurosci.}, vol. 3, no. 1, pp. 71--86, 1991.

\bibitem{b8} P. N. Belhumeur, J. P. Hespanha, and D. J. Kriegman, ``Eigenfaces vs. Fisherfaces: Recognition using class specific linear projection,'' \textit{IEEE Trans. Pattern Anal. Mach. Intell.}, vol. 19, no. 7, pp. 711--720, Jul. 1997.

\bibitem{b9} T. Ahonen, A. Hadid, and M. Pietik\"{a}inen, ``Face description with local binary patterns: Application to face recognition,'' \textit{IEEE Trans. Pattern Anal. Mach. Intell.}, vol. 28, no. 12, pp. 2037--2041, Dec. 2006.

\bibitem{b10} Y. Taigman, M. Yang, M. Ranzato, and L. Wolf, ``DeepFace: Closing the gap to human-level performance in face verification,'' in \textit{Proc. IEEE Conf. Comput. Vis. Pattern Recognit. (CVPR)}, 2014, pp. 1701--1708.

\bibitem{b11} F. Schroff, D. Kalenichenko, and J. Philbin, ``FaceNet: A unified embedding for face recognition and clustering,'' in \textit{Proc. IEEE Conf. Comput. Vis. Pattern Recognit. (CVPR)}, 2015, pp. 815--823.

\bibitem{b12} J. Deng, J. Guo, N. Xue, and S. Zafeiriou, ``ArcFace: Additive angular margin loss for deep face recognition,'' in \textit{Proc. IEEE/CVF Conf. Comput. Vis. Pattern Recognit. (CVPR)}, 2019, pp. 4690--4699.

\bibitem{b13} H. Wang et al., ``CosFace: Large margin cosine loss for deep face recognition,'' in \textit{Proc. IEEE/CVF Conf. Comput. Vis. Pattern Recognit. (CVPR)}, 2018, pp. 5265--5274.

\bibitem{b14} W. Liu, Y. Wen, Z. Yu, M. Li, B. Raj, and L. Song, ``SphereFace: Deep hypersphere embedding for face recognition,'' in \textit{Proc. IEEE Conf. Comput. Vis. Pattern Recognit. (CVPR)}, 2017, pp. 6738--6746.

\bibitem{b15} X. Zhu, Z. Lei, J. Yan, D. Yi, and S. Z. Li, ``High-fidelity pose and expression normalization for face recognition in the wild,'' in \textit{Proc. IEEE Conf. Comput. Vis. Pattern Recognit. (CVPR)}, 2015, pp. 787--796.

\bibitem{b16} Z. Zhu, P. Luo, X. Wang, and X. Tang, ``Multi-view perceptron: A deep model for learning face identity and view representations,'' in \textit{Proc. Adv. Neural Inf. Process. Syst. (NeurIPS)}, 2014, pp. 217--225.

\bibitem{b17} M. Kan, S. Shan, H. Chang, and X. Chen, ``Stacked progressive auto-encoders (SPAE) for face recognition across poses,'' in \textit{Proc. IEEE Conf. Comput. Vis. Pattern Recognit. (CVPR)}, 2014, pp. 1883--1890.

\bibitem{b18} O. M. Parkhi, A. Vedaldi, and A. Zisserman, ``Deep face recognition,'' in \textit{Proc. Brit. Mach. Vis. Conf. (BMVC)}, 2015, pp. 41.1--41.12.

\bibitem{b19} J. Zhao, L. Xiong, P. K. Jayashree, J. Li, F. Zhao, Z. Wang, and S. Yan, ``Dual-agent GANs for photorealistic and identity preserving profile face synthesis,'' in \textit{Proc. Adv. Neural Inf. Process. Syst. (NeurIPS)}, 2017, pp. 66--76.

\bibitem{b20} S. Cole, \textit{Suspect Identities: A History of Fingerprinting and Criminal Identification}. Cambridge, MA, USA: Harvard Univ. Press, 2001.

\bibitem{b21} A. A. Ross, K. Nandakumar, and A. K. Jain, \textit{Handbook of Multibiometrics}. New York, NY, USA: Springer, 2006.

\bibitem{b22} P. Viola and M. Jones, ``Rapid object detection using a boosted cascade of simple features,'' in \textit{Proc. IEEE Conf. Comput. Vis. Pattern Recognit. (CVPR)}, 2001, vol. 1, pp. 511--518.

\bibitem{b23} V. Bazarevsky, Y. Kartynnik, A. Vakunov, K. Raveendran, and M. Grundmann, ``BlazeFace: Sub-millisecond neural face detection on mobile GPUs,'' in \textit{Proc. CVPR Workshop Comput. Vis. Mobile (CVPM)}, 2019.

\bibitem{b24} D. G. Lowe, ``Distinctive image features from scale-invariant keypoints,'' \textit{Int. J. Comput. Vis.}, vol. 60, no. 2, pp. 91--110, Nov. 2004.

\bibitem{b25} K. Zuiderveld, ``Contrast limited adaptive histogram equalization,'' in \textit{Graphics Gems IV}. Cambridge, MA, USA: Academic, 1994, pp. 474--485.

\bibitem{b26} ``ONNX: Open Neural Network Exchange,'' 2019. [Online]. Available: https://onnx.ai

\bibitem{b27} M. Sandler, A. Howard, M. Zhu, A. Zhmoginov, and L.-C. Chen, ``MobileNetV2: Inverted residuals and linear bottlenecks,'' in \textit{Proc. IEEE/CVF Conf. Comput. Vis. Pattern Recognit. (CVPR)}, 2018, pp. 4510--4520.

\end{thebibliography}

\end{document}
